% !TEX TS-program = xelatex
\documentclass[11pt,a4paper]{moderncv}

% --- moderncv theme ---
\moderncvstyle{classic}
\moderncvcolor{blue}

% --- encoding & geometry ---
\usepackage[utf8]{inputenc}
\usepackage[scale=0.8]{geometry}

% --- links ---
% \usepackage[hidelinks]{hyperref}

% --- symbols ---
\usepackage{eurosym} % \euro{} produces €

% --- bibliography ---
\usepackage[backend=biber,style=numeric,sorting=ydnt,defernumbers=true]{biblatex}
\addbibresource{publications.bib}

% ================== Header ==================
\name{Petter}{Törnberg}
\title{Digital Politics / \\ Computational Social Science}
% \title{Computational Social Science}
\address{}{ILLC, University of Amsterdam}
\email{petter.tornberg@uva.nl}
\homepage{pettertornberg.com}
\social[orcid][orcid.org/0000-0001-8722-8646]{0000-0001-8722-8646}
\social[googlescholar][scholar.google.com/citations?hl=en\&user=KuJMuvIAAAAJ]{Google Scholar}
\social[researchgate][www.researchgate.net/profile/Petter-Toernberg]{ResearchGate}
\extrainfo{\href{https://bsky.app/profile/pettertornberg.com}{Bluesky}}

\begin{document}
\makecvtitle

% ================== Profile (tight, metrics-forward) ==================
\section{Profile}
\cvitem{}{Associate Professor in \textit{Computational Social Science} and Head of Research Unit \textit{AI, Culture \& Society} at the \textit{ILLC}, University of Amsterdam. Dutch Research Council (NWO) \textbf{VENI} and \textbf{VIDI} laureate and member of \textit{De Jonge Akademie}. Research at the intersection of digital platforms, AI, and society. Author of 2 books and 70+ publications (e.g., \textit{PNAS}, \textit{New Media \& Society}, \textit{Social Media \& Society}, \textit{Artificial Intelligence Review}); 4{,}300+ citations; 1{,}400 citations in 2025. Frequently featured in outlets such as \textit{The Atlantic}, \textit{The Guardian}, \textit{Ars Technica}, and \textit{Science}.}

% ================== Selected Achievements (new, high-impact hook) ==================
\section{Selected Achievements}
\cvitem{}{%
\begin{itemize}
  \item \textbf{Awarded competitive individual grants:} NWO \textbf{VIDI} (\euro{}850{,}000, comparable to ERC Consolidator) and \textbf{VENI} (\euro{}250{,}000, comparable to ERC StG). Both awarded with a 6–8\% success rate. 
  \item \textbf{Leadership:} Head of \textit{AI, Culture \& Society} research unit (ILLC, UvA); Member, \textit{De Jonge Akademie} (KNAW).
  \item \textbf{Publications:} \textbf{2 books}, \textbf{70+} publications in  venues such as \textit{PNAS}, \textit{New Media \& Society}, \textit{Big Data \& Society}. 
  \item \textbf{Impact:} \textbf{\href{https://scholar.google.com/citations?hl=en&user=KuJMuvIAAAAJ&view_op=list_works}{4{,}300+ citations}}, over 1{,}400 citations in 2025; frequent coverage in major news outlets.
  \item \textbf{Teaching recognition:} Excellent student evaluations (4.7/5). Recipient of the \textit{Collaborative Education Award} (2024); nominated for \textit{Innovative Education Award} (2024). 
\end{itemize}
}

% ================== Major Research Grants (moved up) ==================
\section{Major Research Grants}
\cvitem{2026--2029}{\textbf{Swedish Research Council (VR)} (\euro{}472{,}000). \textit{False words, failing democracies? A global comparative analysis of institutional vulnerability to political misinformation}. \textbf{Co-PI}. Using a global database of social media messages from political parties for a comparative approach to political misinformation. International collaboration with Anton Törnberg and Juliana Chueri.
}
\cvitem{2025--2030}{\textbf{NWO VIDI} (\euro{}850{,}000). \textit{Improving social media with large language models}. \textbf{PI}. (Dutch mid-career grant comparable to an ERC Consolidator). Examining how AI can improve social media for prosocial outcomes. 
}
\cvitem{2024--2027}{\textbf{Trans-Atlantic Platform} (\euro{}500{,}000). \textit{DemDialogueAI -- A Cross-National Study to Promote Democratic Dialogue on Social Media Using Generative Artificial Intelligence}. \textbf{Co-PI}. International collaboration with Chris Bail (Duke), Chris Barrie (NYU), Aniko Hannak (UZH), on simulating social media using LLMs.
}
\cvitem{2021--2025}{\textbf{NWO VENI} (\euro{}250{,}000). \textit{Seeing the city through digital platforms} (VI.Veni.201S.006). \textbf{PI}. (Dutch early-career grant comparable to an ERC Starting Grant.) How digital platforms transform cities and urban conflicts. % TODO: Add one-line impact/outputs (e.g., key pubs, datasets).
}
\cvitem{2017--2021}{\textbf{EU H2020} (\euro{}6{,}000{,}000). \textit{ODYCCEUS} (Grant No.\ 732942). Co-wrote two WPs; funded 4 years postdoc position with 2 PhD positions under my supervision. Focus on how social media reshape political conflicts. % TODO: Add brief role highlight (e.g., WP leadership, coordination tasks).
}
% ================== Academic Appointments ==================
\section{Academic Appointments}
\cventry{2024--}{Associate Professor [UD1], Computational Social Science}{Institute for Logic, Language, and Computation, University of Amsterdam}{Amsterdam, Netherlands}{Tenured. Promotion to UDH2 in process, 2026.}{}
\cventry{2023--2024}{Assistant Professor [UD2], Computational Social Science}{Institute for Logic, Language, and Computation, University of Amsterdam}{Amsterdam, Netherlands}{}{}
\cventry{2022--2025}{Senior Researcher, Geography}{University of Neuchâtel}{Neuchâtel, Switzerland}{Part of \textit{NCCR – On the Move}, a Swiss National Centre of Competence in Research on migration and mobility. Focus on digitalization, governance, and mobility.}{}
\cventry{2021--2025}{Postdoctoral Researcher (VENI), Urban Geography}{University of Amsterdam}{Amsterdam, Netherlands}{Funded through individual NWO VENI grant.}{} 
\cventry{2019--2021}{Postdoctoral Researcher, Complex Systems}{Chalmers University of Technology}{Gothenburg, Sweden}{Mapping online cultural/political conflict. Funded by H2020 ODYCCEUS.}{}
\cventry{2017--2019}{Postdoctoral Researcher, Sociology}{University of Amsterdam}{Amsterdam, Netherlands}{Funded by H2020 ODYCCEUS.}{}

% ================== Education ==================
\section{Education}
\cventry{2022}{Habilitation / P.D. / Associate Professor}{Chalmers University of Technology}{Gothenburg, Sweden}{Title: \textit{Social Media and Polarization: Towards a New Paradigm}}{Referees: Agnieszka Leszczynski, Lasse Gerrits}
\cventry{2017}{PhD in Complex Systems}{Chalmers University of Technology}{Gothenburg, Sweden}{}{Dissertation: ``Worse than Complex'' (Committee Chair: David Byrne; Supervisors: Claes Andersson, David A.\ Lane, Kristian Lindgren)}
\cventry{2011}{MSc in Complex Adaptive Systems}{Chalmers University of Technology}{Gothenburg, Sweden}{}{}
\cventry{2010}{BSc in Computer Science}{Chalmers University of Technology}{Gothenburg, Sweden}{}{}

% ================== PhD Supervision ==================
\section{PhD Supervision}
\cvitem{2025--}{\textbf{Erfan Monazzah} (ILLC, UvA): Improving social media through AI. Funded by VIDI grant.
}
\cvitem{2025--}{\textbf{Michelle Schimmel} (ILLC, UvA): Language models, social media, and affective polarization. Funded by VIDI grant.
}
\cvitem{2025--}{\textbf{Daniele Barolo} (Complexity Science Hub, Vienna): \textit{Advisory team.} Large language models as tools for social science research.
}
\cvitem{2018--2025}{\textbf{Anna Keuchenius} (Sociology, UvA): ``A computational social science that makes sense: Studying the complex dynamics of social interaction and collective meaning-making.'' Funded by EU H2020 ODYCCEUS.}
\cvitem{2018--2023}{\textbf{Livia van Vliet} (Sociology, UvA): ``Digital soapboxes: Analyzing Twitter politics across 26 countries.'' Funded by EU H2020 ODYCCEUS.}

% ================== MSc and BSc Supervision ==================
\section{Recent MSc and BSc Supervision}
\cvitem{2025--2026}{\textbf{Xuening Tang} (MSc, UvA): Beyond Black-Box Agents: Explainable and Validatable Generative ABMs (research paper in draft.)
}
\cvitem{2025--2025}{\textbf{Mark Larooij} (MSc, UvA): Simulating social media with LLM (thesis resulted in two research papers; see below.)
}
\cvitem{2025--2025}{\textbf{Victor Hartman} (MSc, UvA): Measuring negative campaigning with LLMs (thesis resulted in one research paper; see below.)
}
\cvitem{2025--2025}{\textbf{Kyriaki Polykarpou} (MSc, UvA): Simulating viral misinformation with Generative ABMs.}
\cvitem{2024--2025}{\textbf{Ilse Nederveen} (BSc, UvA): Gender and racial biases in LLMs.}

% ================== Teaching (front-load recognition) ==================
\section{Teaching}
\cvitem{}{Extensive experience designing and teaching courses that combine critical and applied perspectives on digital platforms, political communication, and computational methods. Courses consistently receive excellent evaluations (4.7/5.0). Recipient of the \textit{Collaborative Education Award} (2024) and nominee for the \textit{Innovative Education Award} (2024).}
\cventry{2023--}{MSc Course: Computational Social Science Analysis}{Social Science Research Master, University of Amsterdam}{Methods and theories for studying a digital society.}{Developer, coordinator, lecturer. \textbf{Nominated} for \textit{Innovative Education Award} (2024).}{}
\cventry{2023--}{BSc Program: Computational Social Science}{University of Amsterdam}{Exploring how data and algorithms shape society, politics, and communication.}{Lecturer and program co-developer. \textbf{Winner} \textit{Collaborative Education Award} (2024).}{}
\cventry{2014--2016}{Simulation of Complex Systems}{Chalmers University of Technology}{}{Course lecturer and MSc supervisor.}{}
\cventry{2014--2016}{Complex Systems Seminars}{Chalmers University of Technology}{}{Course coordinator and lecturer (MSc).}{}
\cventry{2014}{Social Network Analysis}{Chalmers University of Technology}{}{Coordinator, lecturer, and course creator (PhD level).}{}
% ================== Leadership & Academic Service ==================
\section{Leadership \& Academic Service}
\cvitem{}{Demonstrated leadership in building research units, shaping curricula, and contributing to national and international academic communities.}
\cventry{2024--}{Head of Research Unit}{AI, Culture \& Society, ILLC, UvA}{}{}{}
\cventry{2024--}{Member of \textit{De Jonge Akademie}}{ (KNAW; Royal Netherlands Academy of Arts and Sciences)}{ a platform recognizing outstanding early-career researchers across disciplines}{}{} % kept concise; long description moved to achievements
\cventry{2024--}{EU Horizon External Expert}{Reviewer for Horizon Grants on Disinformation}{}{}{}
\cventry{2023--}{Examination Committee Member}{MoL Examination Committee, ILLC, UvA}{}{}{}
\cventry{2023--}{BSc Program Developer}{CCSci, University of Amsterdam}{}{Part of the founding team of the award-winning program at the intersection of technology and society.}{}
\cventry{2020--}{Peer reviewer}{\textit{Science}, \textit{PNAS}, \textit{Nature Human Behavior}, \textit{Nature Communications}, \textit{New Media \& Society}, \textit{etc.}}{}{}{}
\cventry{2019--2021}{Committee Member}{Program \& Workshop committees: ANNSIM HSAA}{}{}{}

% ================== Recent Keynotes ==================
\section{Recent Keynotes \& Invited Talks}
\cventry{2025}{``Studying Digital Media in the \textit{Post-Social} Media Era''}{The Future of Social Media Research, Oxford}{Keynote}{}{}
\cventry{2025}{``How social media shape political discourse across 200 countries''}{Networks, (mis)information \& Polarization Conference, Utrecht University}{Keynote}{}{}
\cventry{2025}{``Can We Fix Social Media?''}{Michael Bernstein Lab, Stanford University}{Invited Talk}{}{}
\cventry{2025}{``Large Language Models and Social Simulation''}{Lab42 Conference, UvA}{Keynote}{}{}
\cventry{2024}{``Improving Social Media with LLMs''}{Machine Learning in Science Conference, University of Glasgow}{Keynote}{}{}
\cventry{2024}{``Social Media and Polarization: Towards a New Paradigm''}{Symposium on Social Cohesion, Hamburg}{Keynote}{}{} 
\cventry{2024}{``Validation and Calibration in Social Simulation''}{University of Oxford}{Invited Talk}{}{}
\cventry{2023}{``Simulating Social Media Platforms with Large Language Models''}{Generative Methods conference, Aalborg University, Copenhagen}{Keynote}{}{} % 

% ================== Professional Training ==================
\section{Professional Training}
\cventry{2023}{UTQ/BKO Dutch Teaching Certification}{University of Amsterdam}{}{}{}
\cventry{2021}{Veni/Vidi/Vici ERC Leadership in Research}{University of Amsterdam}{}{}{}
\cventry{2019}{Supervising Research Students (CLS905)}{Chalmers University of Technology}{}{}{}
\cventry{2013}{Teaching, Learning \& Evaluation (FCIU010)}{Chalmers University of Technology}{}{}{}

% ================== Languages ==================
\section{Languages}
\cvitem{}{Swedish (Native); English (Fluent); Portuguese (B2); Dutch (B1); German (B1); Spanish (B1); Italian (A2); French (A2).}

\newpage 
% ================== Publications ==================
\section{Publications}
\cvitem{}{\textbf{\href{http://scholar.google.com/citations?hl=en&user=KuJMuvIAAAAJ&view_op=list_works}{See Google Scholar}.} Citations: 4{,}310; h-index: 29; i10-index: 46. % TODO: Update metrics as they change.
}

% --- Journalistic work (kept concise) ---
\subsection*{Journalistic Work}
\cventry{2025}{``The Ascendance of Algorithmic Tyranny''}{\textit{Noema}}{\url{https://www.noemamag.com/the-ascendance-of-algorithmic-tyranny}}{}{}
\cventry{2023}{``Visst är AI komplexa system -- men det finns större problem''}{\textit{Dagens Nyheter}}{\url{https://www.dn.se/kultur/petter-tornberg-visst-ar-ai-komplexa-system-men-det-det-finns-storre-problem/}}{}{}
\cventry{2021}{``De online echoput av extreem-rechts''}{\textit{De Groene Amsterdammer}}{\url{https://www.groene.nl/artikel/de-online-echoput-van-extreem-rechts}}{}{}

% --- Books ---
\subsection*{Books}
\nocite{*}
\printbibliography[type=book,heading=none,resetnumbers=true]

% --- Optional: Selected publications block (add keyword={selected} in .bib) ---
% TODO: If you want a short, skimmable list, tag 8–12 entries in publications.bib with keyword={selected}
% and uncomment the block below to show them first.
% \subsection*{Selected Peer-Reviewed Articles}
% \printbibliography[type=article,keyword=selected,heading=none,resetnumbers=true]

\subsection*{Forthcoming}
\begin{itemize}
    \item Törnberg P. It's not Polarization; It's the Radicalization of the Political Right. \textit{Perspectives on Politics} (Major Revision)
    \item Törnberg P. J Chueri. Elite Political Discourse has Become More Toxic in Western Countries. \textit{Nature Human Behavior} (Major Revision)
    \item Larooij K, P Törnberg. Can We Fix Social Media? Testing Prosocial Interventions using Generative Social Simulation. \textit{Nature Communications} (R\&R)
    \item Pagan N, P Törnberg, CA Bail, A Hannák, C Barrie. Computational Turing Test Reveals Systematic Differences Between Human and AI Language \textit{PNAS}  (Major Revision)
    \item Törnberg, P. Shifts in U.S. Social Media Use, 2020–2024: Decline, Fragmentation, and Enduring Polarization. \textit{Journal of Quantitative Description: Digital Media} (Major Revision)
    \item Hartman V. Törnberg, P. Who attacks, and why? Using LLMs to identify negative campaigning in 18M tweets across 19 countries. \textit{Political Communication} (Major Revision)
    \item Törnberg, P. The City as an Anti-Growth Machine. \textit{Antipode} (Major Revision)
    
\end{itemize}


% --- Full list of journal articles ---
\subsection*{Peer-Reviewed Journal Articles}
\printbibliography[type=article,heading=none,resetnumbers=true]

% --- Chapters, preprints, proceedings (condensed) ---
\subsection*{Chapters, Preprints, and Conference Proceedings}
\printbibliography[nottype=article,nottype=book,heading=none,resetnumbers=false]

\end{document}

% % !TEX TS-program = xelatex
% \documentclass[11pt,a4paper]{moderncv}
% % moderncv themes
% \moderncvstyle{classic}
% \moderncvcolor{blue}
% \usepackage[utf8]{inputenc}
% \usepackage[scale=0.8]{geometry}
% % Bibliography management
% \usepackage[backend=biber,style=numeric,sorting=ydnt,defernumbers=true]{biblatex}
% \usepackage{eurosym}
% \addbibresource{publications.bib}% BibTeX file with all publications

% \name{Petter}{Törnberg}
% \title{Associate Professor in Computational Social Science}
% \address{}{ILLC, University of Amsterdam}
% \email{petter.tornberg@uva.nl}
% \homepage{pettertornberg.com}
% \social[orcid][https://orcid.org/0000-0001-8722-8646]{0000-0001-8722-8646}
% \social[researchgate][https://www.researchgate.net/profile/Petter-Toernberg]{ResearchGate}
% \extrainfo{\href{https://bsky.app/profile/pettertornberg.com}{Bluesky}}


% \begin{document}
% \makecvtitle

% %---- Profile ----
% \section{Profile}
% \cvitem{}{Associate Professor in \textit{Computational Social Science} and Head of \textit{AI, Culture \& Society} at the \textit{ILLC}, University of Amsterdam. NWO VENI and VIDI laureate and member of \textit{De Jonge Academie}. His interdisciplinary research focuses on the intersection of digital platforms, AI, and society. He has published 2 books, and over 60 papers and chapter, in outlets such as \textit{PNAS}, \textit{New Media \& Society}, and \textit{Big Data \& Society}. He is frequently featured in media outlets such as \textit{The Atlantic}, \textit{The Guardian}, \emph{New Scientist}, and \textit{Insider}.}

% % \cvitem{}{Interdisciplinary scholar focused on how platforms and algorithms intersect with social and political life. Combines large-scale text and network analysis with LLM-based social simulation. Author of two books and 60+ publications (e.g., \textit{PNAS}, \textit{New Media \& Society}, \textit{Big Data \& Society}); frequent keynote speaker and media commentator.


% %---- Education ----
% \section{Education}
% \cventry{2022}{Habilitation / P.D. / Associate Professor}{Chalmers University of Technology}{Gothenburg, Sweden}{}{Referees: Agnieszka Leszczynski, Lasse Gerrits}
% \cventry{2017}{PhD in Complex Systems}{Chalmers University of Technology}{Gothenburg, Sweden}{}{Dissertation: "Worse than Complex" (Committee Chair: David Byrne; Supervisors: Claes Andersson, David A. Lane, Kristian Lindgren)}
% \cventry{2011}{MSc in Complex Adaptive Systems}{Chalmers University of Technology}{Gothenburg, Sweden}{}{ }
% \cventry{2010}{BSc in Computer Science}{Chalmers University of Technology}{Gothenburg, Sweden}{}{ }

% %---- Academic Positions ----
% \section{Academic Appointments}
% \cventry{2024--}{Associate Professor [UD1], Computational Social Science}{Institute for Logic, Language, and Computation, University of Amsterdam}{Amsterdam, Netherlands}{Tenured position. Promotion to UDH2 January 2026}{}
% \cventry{2023--2024}{Assistant Professor [UD2], Computational Social Science}{Institute for Logic, Language, and Computation, University of Amsterdam}{Amsterdam, Netherlands}{}{}
% \cventry{2022--2025}{Senior Researcher, Geography}{University of Neuchâtel}{Neuchâtel, Switzerland}{Part of \textit{NCCR - On the Move}}{ }
% \cventry{2021--2025}{NWO VENI Postdoctoral Fellow, Urban Geography}{University of Amsterdam}{Amsterdam, Netherlands}{Funded through individual NWO VENI grant.}{}
% \cventry{2019--2021}{Postdoctoral Researcher, Complex Systems}{Chalmers University of Technology}{Gothenburg, Sweden}{Mapping online cultural/political conflict. Funded through Horizon2020 grant ODYCCEUS.}{}
% \cventry{2017--2019}{Postdoctoral Researcher, Sociology}{University of Amsterdam}{Amsterdam, Netherlands}{Funded through Horizon2020 grant ODYCCEUS.}{}

% %---- Grants ----
% \section{Major Research Grants}
% \cvitem{2025--2030}{\textbf{NWO VIDI} (\euro{}850k). ``Improving social media with large language models''. Individual, highly competitive mid-career grant. Similar to \textit{ERC Consolidator}.}
% \cvitem{2021--2025}{\textbf{NWO VENI} (\euro{}250k). ``Seeing the city through digital platforms'' (VI.Veni.201S.006). Individual, highly competitive early-career grant, similar to \textit{ERC Starting Grant}.}
% \cvitem{2017--2021}{\textbf{EU H2020} (\euro{6M}). ``ODYCCEUS'' (Grant No.\ 732942). I co-wrote two WPs, which funded by my postdoc position, and 2 PhD students under my supervision. }

% %---- Supervision ----
% \section{PhD Supervision}

% \cvitem{2025--}{\textbf{Michelle Schimmel} (ILLC): Language models, social media, and affective polarization.}{}

% \cvitem{2018--2025}{\textbf{Anna Keuchenius} (Sociology): ``A computational social science that makes sense: Studying the complex dynamics of social interaction and collective meaning-making.''}{}

% \cvitem{2018--2023}{\textbf{Livia van Vliet} (Sociology): ``Digital soapboxes: Analyzing Twitter politics across 26 countries.''}{}

% %---- Teaching ----
% \section{Teaching}
% \cventry{2023--}{MSc course: Computational Social Science Analysis}{Social Science Research Masters, University of Amsterdam}{}{Developer, coordinator, lecturer. \textit{Nominated for 2024 Innovative Education Award}}{ }
% \cventry{2023--}{BSc Computational Social Science}{University of Amsterdam}{}{Lecturer and program co-developer. \textit{Winner of 2024 Collaborative Education Award}}{ }
% \cventry{2014--2016}{Simulation of Complex Systems}{Chalmers University of Technology}{}{Course lecturer and supervisor at MSc level}{ }
% \cventry{2014--2016}{Complex Systems Seminars}{Chalmers University of Technology}{}{Course coordinator and lecturer at MSc level}{ }
% \cventry{2014}{Social Network Analysis}{Chalmers University of Technology}{}{Coordinator, lecturer, and course creator at PhD level}{ }

% %---- Service ----
% \section{Leadership \& Academic Service}
% \cventry{2024--}{Head of Research Unit}{`AI, Culture \& Society', ILLC, UvA}{}{}{}
% \cventry{2024--}{Member of \textit{De Jonge Akademie}}{\textit{De Jonge Akademie} is a platform within the Royal Netherlands Academy of Arts and Sciences (KNAW) for outstanding early- to mid-career scholars from across disciplines. Members are selected for their academic excellence and engagement with broader societal and interdisciplinary issues.}{}{}{}
% \cventry{2023--}{Examination Committee Member}{Member of the MoL examination committee, ILLC, UvA}{}{}{}
% \cventry{2023--}{BSc Program Developer}{Part of team developing the award-winning BSc program CCSci at University of Amsterdam.}{}{}{}
% \cventry{2020--}{Peer-reviewer}{Frequent reviewer for: \emph{Science}, \emph{PNAS}, \emph{Nature Human Behavior}, \emph{Nature Communication}, \emph{New Media \& Society}, etc.}{}{}{}
% \cventry{2019--2021}{Committee Member}{Program \& Workshop committees: ANNSIM HSAA}{}{}{}

% %---- Invited Talks ----
% \section{Recent Keynotes}
% \cventry{2025}{``How social media shape political discourse across 260 countries''}{Networks, (mis)information \& polarization conference, Utrecht University}{}{}{ }
% % \cventry{2025}{"Global comparative perspectives on social media and polarization"}{Michael Bernstein Group, Stanford University}{}{}{ }
% \cventry{2024}{``Large Language Models and Text Analysis''}{Lab 42 Conference, UvA}{}{}{ }
% \cventry{2024}{``Improving Social Media with LLMs''}{Machine Learning in Science Conference, University of Glasgow}{}{}{ }
% \cventry{2024}{``Social media and Polarization: Towards a new paradigm''}{Symposium on Social Cohesion, Hamburg}{}{}{ }

% %---- Professional Training ----
% \section{Professional Training}
% \cventry{2023}{UTQ/BKO Dutch Teaching Certification}{University of Amsterdam}{}{}{ }
% \cventry{2021}{Veni/Vidi/Vici ERC Leadership in Research}{University of Amsterdam}{}{}{ }
% \cventry{2019}{Supervising Research Students (CLS905)}{Chalmers University of Technology}{}{}{ }
% \cventry{2013}{Teaching, Learning \& Evaluation (FCIU010)}{Chalmers University of Technology}{}{}{ }


% %---- Languages ----
% \section{Languages}
% \cvitem{}{Swedish (Native); English (Fluent); Portuguese (B2); Dutch (B1); German (B1);  Spanish (B1); French (A2)}

% %---- Publications ----

% %---- Metrics ----
% \section{Publications}

% \textbf{\href{http://scholar.google.com/citations?hl=en&user=KuJMuvIAAAAJ&view_op=list_works}{See \underline{Google Scholar}.} Citations: 3502; H-index: 27; i10-index: 42.}

% % Include all entries
% \nocite{*}

% \subsection{Journalistic Work}
% \cventry{2025}{``The Ascendance of Algorithmic Tyranny''}{Noema}{\url{https://www.noemamag.com/the-ascendance-of-algorithmic-tyranny}}{}{ }
% \cventry{2023}{``Visst är AI komplexa system -- men det finns större problem''}{(`AI are indeed complex systems -- but there are more fundamental risks'') Dagens Nyheter}{\url{https://www.dn.se/kultur/petter-tornberg-visst-ar-ai-komplexa-system-men-det-det-finns-storre-problem/}}{}{ }
% \cventry{2021}{``De online echoput av extreem-rechts''}{(``The online echo chamber of the extreme right'') De Groene Amsterdammer}{\url{https://www.groene.nl/artikel/de-online-echoput-van-extreem-rechts}}{}{ }

% \subsection*{Books}
% \printbibliography[type=book,heading=none,resetnumbers=true]

% \subsection*{Chapters, Preprints, and Conference Proceedings}
% % include remaining types: proceedings, misc, etc.
% \printbibliography[nottype=article,nottype=book,heading=none,resetnumbers=true]


% \subsection*{Peer-Reviewed Journal Articles}
% \printbibliography[type=article,heading=none,resetnumbers=true]



% % \subsection*{Books}
% % \nocite{*}
% % \printbibliography[type=book,heading=none]

% % % The following command prints entries from the .bib file tagged with 'article'
% % \subsection*{Selected Articles}
% % \nocite{*}
% % \printbibliography[heading=none]
% % %[type=article,heading=none]

% \end{document}


% 12-09-2025
%  Interview iwth science margainze 
% Hannah Richter | she/her
% Science Journalist
% Outreach Coordinator - Next Generation Event Horizon Telescope
% 610-635-9807
% hannah-richter.com
% Washington, D.C. (ET)


% 12-09-2025
% Jennifer Ouellette <jennifer.ouellette@arstechnica.com>
% Interview with Ars Technica about Larooij and Tornberg 2025.

% TechLinked 1.98M subscri 
% One of the world's largest Tech YOutube 
% https://www.youtube.com/watch?v=jJR_Mt50gSI&t=485s

% Futurism
% https://futurism.com/social-network-ai-intervention-echo-chamber 


% Lisa Peters 
% Stichting Democratie en Media (SDM)
% NGO working with democracy and media, 


% 23-09-2025
% Sean Carrol's Mindscape Podcast Interview


% From: Haug, Clemens <Clemens.Haug@mdr.de>
% Date: Thursday, 21 August 2025 at 11:44
% To: Petter Törnberg <p.tornberg@uva.nl>
% Subject: Media // Can We Fix Social Media?
 % Mitteldeutscher Rundfunk, Leipzig, Germany

 % Interview with Justin Hendrix
 % 2x 

 % Kurzgezagt YouTube channel with 25.6M viewers made a video of my PNAS paper. 
 